\section{Introduction}

\subsection{Motivation}

The softmax attention operation computes a pairwise score between every 
query and every key:

\[
\text{Attn}(Q, K, V) = \text{softmax}\!\left(\frac{Q K^\top}{\sqrt{d_h}}\right) V
\]

This requires $\mathcal{O}(L^2 d)$ time and $\mathcal{O}(L^2)$ memory for 
sequence length $L$. For $L = 4096$ (a common long-context setting), the 
attention matrix alone requires 64 MB per head, per layer, in FP16.

\textbf{Linear attention} reduces this to $\mathcal{O}(L d)$ or 
$\mathcal{O}(L \log L)$ by either restricting attention (sparse, local 
windows) or approximating the softmax kernel. Both approaches come with 
trade-offs:

\begin{itemize}
\item \textbf{Sparse attention} (Longformer, BigBird) captures local structure well
  but loses global context.
\item \textbf{Kernel attention} (Performer, Linear Transformer) preserves global
  context but approximates softmax, which is problematic when the true
  attention is peaked.

\end{itemize}
\textbf{Cluster-based attention} offers a third path. Instead of comparing 
individual tokens, group them into $K$ clusters, compute attention at the 
\textbf{cluster level}, and redistribute the results. If $K \ll L$, this is much 
cheaper than dense attention while preserving \textbf{global reach} (every 
cluster attends to every key).

\subsection{Prior Work on Clustering Attention}

\textbf{Routing Transformer} (Roy et al., 2021) clusters tokens and restricts 
attention within each cluster. Its complexity is $\mathcal{O}(L \cdot W)$ 
but requires a large cluster size $W$ to maintain context, negating the 
savings.

\textbf{Reformer} (Kitaev et al., 2020) uses locality-sensitive hashing (LSH) 
to approximate attention. Its complexity depends on hash collisions and is 
not guaranteed to be linear.

\textbf{Cluster-Former} (Wang et al., 2021) uses clustering for document-level 
tasks but applies it at a coarse granularity (sentence-level).

\textbf{ClusterFormer} differs by computing \textbf{centroid-to-key} attention 
directly, giving exact global context at $\mathcal{O}(K L d)$ complexity.

\subsection{Problem Statement}

Given a pretrained attention module with projections $W_Q, W_K, W_V, W_O$, 
we seek a modified attention function $\mathcal{A}'$ such that:

\begin{enumerate}
\item \textbf{Subquadratic complexity.} $\mathcal{A}'$ runs in $\mathcal{O}(K L d)$
  with $K \ll L$.
\item \textbf{Global context.} Every query attends to a representation of the whole
  sequence, not just a window.
\item \textbf{Local structure.} Nearby syntactic dependencies are preserved.
\item \textbf{Step-0 preservation.} With zero-init LoRA adapters, the initial forward
  pass matches the pretrained model.
\item \textbf{Multi-architecture support.} Works on flat, nested, and causal
  transformers.

\end{enumerate}
\subsection{Our Approach}

ClusterFormer computes attention in three steps:

\textbf{Step 1: Cluster queries.} Apply fast differentiable K-means to the query 
representations $Q \in \mathbb{R}^{B \times H \times L \times d_h}$, 
producing $K$ centroids $C \in \mathbb{R}^{B \times H \times K \times d_h}$.

\textbf{Step 2: Centroid-to-key attention.} Compute attention from each centroid 
to all keys:

\[
V_{\text{cent}} = \text{softmax}\!\left(\frac{C K^\top}{\sqrt{d_h}}\right) V \in \mathbb{R}^{B \times H \times K \times d_h}
\]

This costs $\mathcal{O}(K L d_h)$ -- linear in $L$.

\textbf{Step 3: Redistribute.} Assign each query to its nearest centroid and 
collect the corresponding context:

\[
\text{Global}(Q) = \text{softmax}\!\left(\frac{Q C^\top}{\sqrt{d_h} \cdot \tau}\right) V_{\text{cent}}
\]

where $\tau$ is a learnable temperature.

To capture local structure missed by clustering, we add a sliding-window 
attention branch and blend via a learnable gate $\alpha$:

\[
\text{Output} = (1 - \alpha) \cdot \text{Global}(Q) + \alpha \cdot \text{Local}(Q)
\]

\subsection{Contributions}

\begin{itemize}
\item \textbf{Centroid-to-key cross-attention.} A novel formulation of cluster
  attention that computes exact global context at linear cost. Section 4.3.
\item \textbf{Fast differentiable K-means.} A vectorized Lloyd algorithm with
  padding-aware initialization and 3 iterations, sufficient for attention
  clustering. Section 4.2.
\item \textbf{Learnable temperature and blend.} A temperature parameter $\tau$ for
  assignment sharpening and a blend scalar $\alpha$ for global-local fusion,
  both learned during recovery. Section 4.4.
\item \textbf{Broad architecture support.} Multi-architecture weight extraction plus
  GQA, RoPE, nested encoder bypass, and causal masking. Section 5.
\item \textbf{Systematic evaluation.} Three models, SST-2, ablations over $K$,
  window size, and blend initialization. Section 7.
\item \textbf{Honest disclosure.} Parameter count increases due to LoRA and the
  learnable scalars; we report this transparently. Section 7.2.
\item \textbf{Full reproducibility.} Code, checkpoints, and evaluation scripts
  released under MIT.

\end{itemize}
\subsection{Paper Organization}

Section 2 reviews related work. Section 3 provides background on attention, 
K-means, and linear attention. Section 4 presents ClusterFormer. Section 5 
describes architecture integration. Section 6 details experimental setup. 
Section 7 reports results. Section 8 analyzes complexity, clustering quality, 
and ablation contributions. Section 9 discusses usage. Section 10 lists 
limitations. Section 11 concludes.



\subsection{Notation}

\begin{lstlisting}

---
\end{lstlisting}